\snote{single}

The method requires at least $k_{\min}=2$ biological parameters. This note examines whether a single parameter could suffice.

\subsection*{Design}
For every admissible problem the candidates are the ten biological parameters with the largest sensitivity energy $E_j$ and the ten with the largest finite effect $s_j$ (the two orderings of the method). Each candidate $j$ is tested alone, with the calibration parameters $\Theta_C$ re-estimated together with it as in the method: $\{j\}$ is admissible if the median $e_{\mathrm{fit}}$ over the same 15 scenarios at $\pm5\%$ is at most $\tau_e$ (without the requirement $|S|\ge2$). For the best admissible candidate $j^\ast$ (smallest median $e_{\mathrm{fit}}$) we also computed the median $e_{\mathrm{fit}}$ at $\pm10\%$ and its external validation, real and synthetic, exactly as for the selected subset $S$ (Supplementary Note~\emph{Synthetic prediction experiment}). Bachmann, Chen, Isensee, Lang, Raimundez were not included because the repeated nonlinear re-estimations are too slow.

\subsection*{Results}
\textit{Being computed for 26 problems; this section is generated automatically when complete.}
